\documentclass[14pt,a4paper]{extarticle}
\usepackage[left=3cm,right=2cm,top=2cm,bottom=2cm,headheight=18pt]{geometry}
\usepackage{fontspec}
\setmainfont{Liberation Serif}
\setsansfont{Liberation Serif}
\setmonofont{Liberation Mono}
\usepackage{unicode-math}
\setmathfont{texgyretermes-math.otf}
\usepackage{amsmath,graphicx,booktabs,longtable,array,enumitem}
\usepackage{titlesec,fancyhdr}
\usepackage[hidelinks,unicode]{hyperref}
\usepackage{url}
\urlstyle{same}
\setlength{\parindent}{1.27cm}
\setlength{\parskip}{0pt}
\linespread{1}
\setlength{\emergencystretch}{2em}
\clubpenalty=10000
\widowpenalty=10000
\titleformat{\section}{\normalfont\bfseries}{\thesection.}{0.5em}{}
\titleformat{\subsection}{\normalfont\bfseries}{\thesubsection.}{0.5em}{}
\titlespacing*{\section}{0pt}{1.2em}{0.5em}
\titlespacing*{\subsection}{0pt}{0.9em}{0.3em}
\setlist{nosep,leftmargin=1.27cm}
\pagestyle{fancy}
\fancyhf{}
\renewcommand{\headrulewidth}{0pt}
\renewcommand{\footrulewidth}{0pt}
\fancyhead[C]{\ifnum\value{page}>1\thepage\fi}
% Liberation Serif is a disclosed Times-compatible substitute, not Times New Roman.
% Current statutory geometry/size baseline: O‘RQ-682 appendix §§43–51.
\hyphenpenalty=10000
\exhyphenpenalty=10000
\pretolerance=10000
\tolerance=10000
\setlength{\emergencystretch}{4em}
\renewcommand{\normalsize}{\fontsize{14bp}{16.8bp}\selectfont}
\normalsize
\begin{document}
\noindent Автор: Abduxoliq Akmal ugli Ashuraliyev. Редакция от 1 октября 2026 года.

\begin{flushright}Проект\end{flushright}
\begin{center}\textbf{ПРЕДЛОЖЕНИЕ ДЛЯ ВКЛЮЧЕНИЯ\ В ПРОЕКТ НОРМАТИВНО-ПРАВОВОГО АКТА}\par
\medskip\textbf{О порядке проверки расчётов многомерной оценки семей}\end{center}
Авторское предложение по пункту 14 PF-193. Основная русская редакция авторского предложения для подачи. Согласованные узбекская и английская редакции приложены; узбекская предназначена для работы над официальным проектом на государственном языке. Вид окончательного акта и реквизиты утверждения определяет разработчик. Новые положения ниже предлагаются, а не объявляются действующим правом.

1. Утвердить согласно приложению порядок проверки расчётов методики многомерной оценки, применяемой при включении семей в Социальный реестр, распределении по категориям и направлении к соответствующим услугам.

2. Национальному агентству социальной защиты при Президенте Республики Узбекистан (далее — Агентство) совместно с Министерством экономики и финансов Республики Узбекистан (далее — Министерство) при подготовке методики для утверждения в установленном порядке отразить в едином описании критерии, правила расчёта, границы категорий и их правовое и научное обоснование. Не допускать самостоятельного введения программными настройками нового критерия или основания отказа.

3. Агентству обеспечить проверку соответствия расчётов утверждённой методике, возможность повторения и сохранение оснований решения. Проверка не является основанием продления установленных законодательством сроков рассмотрения обращений семей.

4. Агентству совместно с Министерством организовать подготовку методики и проверочных материалов в порядке статьи 22 Закона «О нормативно-правовых актах». В рабочую группу могут включаться с их согласия граждане, обладающие необходимыми знаниями и опытом; могут запрашиваться научные рекомендации и консультации специалистов. Агентству закрепить методические, системные и юридические обязанности и утвердить внутренний технический регламент паспортов критериев, сопоставления источников, реестра тестов, протоколов и записей дефектов. Регламент исполняет настоящий Порядок и не вводит критерии нуждаемости, основания отказа, обязанности заявителя или изменения законных сроков. Участие в подготовке не даёт полномочия утверждать критерии или получать защищённые записи.

5. Агентству и Министерству согласовать описание расчёта с правилами оценки в реестре, определения категории, помощи, жалоб и обмена информацией в основном проекте до первого применения затронутого правила. Противоречие вышестоящего акта вносится компетентному принимающему органу. Техническое исправление исполняет правило, законно применимое на дату решения; оно не продлевает заменённые критерии, не меняет дату внедрения и не даёт работнику права вручную внести решение о категории вне законной процедуры.

6. Проверку включить в существующий процесс разработки системы, технического согласования и управления изменениями. Использовать действующие спецификации, тесты, записи и законный контроль при наличии эквивалентного доказательства; не требовать повторных контролей и закупок. До обязательства дополнительных расходов Агентству и Министерству определить дополнительные работы, назначенное время работников, покрытие договором и законное финансирование. Изменение обработки обосновывается относящимися к нему записями нагрузки и времени, при необходимости — испытанием мощности в разрешённой среде. Формула стоимости, штатное расписание и общая сумма бюджета не заменяют таких доказательств. Порядок не учреждает отдельную общереспубликанскую службу ручного рассмотрения, параллельную систему, плату с заявителя или дополнительный этап заявления.

7. До юридически значимого применения Агентству зафиксировать законное правило, требуемые технические согласования, результаты контроля и готовность существующих каналов объяснения и возражения в записи выпуска по пункту 25. Неуспешная проверка рассматривается по пункту 26. Эмпирическое исследование благосостояния и отдельный публичный отчёт не являются условиями внедрения. Техническая приёмка не заменяет законные согласования, не меняет начало 1 января 2027 года и не продлевает сроки заявителя.

\clearpage\setcounter{page}{1}
\begin{flushright}\fontsize{12bp}{14.4bp}\selectfont Приложение к материалу 1\\Авторское предложение\\Редакция от 1 октября 2026 года\end{flushright}
\begin{center}\textbf{Порядок проверки расчётов многомерной оценки}\end{center}
\subsection*{Глава 1. Общие положения}
1. Порядок касается соответствия расчётов утверждённой методике, прослеживаемости использованных данных и сохранения оснований решения. Контрольные случаи проверяются до внедрения и относящихся к ним изменений в существующем процессе выпуска, без второго ручного расчёта каждого заявления. Исследования чувствительности или благосостояния требуют отдельно согласованных вопроса, законных данных и ресурсов; они не являются дополнительными условиями помощи или запуска утверждённого правила. Направление к услугам проверяется в части, вытекающей из этого правила; отраслевые условия услуг сохраняются.

2. В настоящем Порядке используются сокращённые наименования: информационная система «Yagona milliy ijtimoiy himoya» (далее — YAMIH) и центры социальных услуг «Инсон» Агентства (далее — центры «Инсон»). Основные понятия:

\textbf{описание методики} — документ о критериях оценки, источниках, правилах расчёта, условиях категорий и сроках применения;

\textbf{версия расчёта} — связанная запись идентификатора утверждённой методики, идентификатора программного выпуска и конфигурации параметров с датами действия, использованных в расчёте;

\textbf{протокол проверки} — запись условий, результатов, несоответствий и мер;

\textbf{неразрешённые данные} — отсутствующие, неподтверждённые, противоречивые либо не отвечающие требованиям обновления методики;

\textbf{контрольный расчёт} — результат независимого повторения по утверждённым правилам.

3. Методика не заменяет полномочие назначать помощь или услугу. Категория, право на конкретную помощь и её размер определяются отдельно по законодательству.

\subsection*{Глава 2. Описание и версии}
4. Для каждого критерия описание определяет правовое основание, цель, уполномоченный источник, единицу, учитываемый период, требования обновления и правило расчёта. Даты применения, меняющиеся параметры, повторяющиеся либо пересекающиеся сведения и исключения определяются в соответствии с утверждённой методикой. Внутренний технический регламент Агентства устанавливает поля паспортов, сопоставление источников и проверки точности; формы не устанавливают самостоятельного нового правила или основания отказа.

5. Описание охватывает формулы, условия, коэффициенты, зависимости, исключения, округление и границы категорий. Указывается категория при точном равенстве. После применения законных исключений и приоритетов проверяется отсутствие противоречивых категорий и неразрешённого промежутка между включением и отказом. Уточнение не смешивается с решением о включении или отказе. Программа не применяет отсутствующее в описании правило.

6. Для каждой версии сохраняются основание утверждения, начало и конец действия, различия с прежней и протокол. Изменение правового последствия утверждается компетентным органом в установленном порядке. Правовая методика, программный выпуск и конфигурация параметров имеют отдельно различимые записи, связанные с версией расчёта; даты правового действия отличаются от дат технического внедрения. Программное исправление исполнения неизменного законного правила не объявляется новым критерием политики. Конфигурация определяет сопоставление источников, преобразования и округление; изменения, способные повлиять на результаты, проверяются.

7. Электронное описание указывает опубликованный акт утверждения методики, даты её применения и версию расчёта. Открытые сведения о правиле доступны через существующий официальный ресурс Агентства; допустима ссылка на опубликованное уполномоченное описание вместо его копирования во второй отчёт. Ограниченная часть доступна только уполномоченным проверяющим и рассматривающим жалобу. Публикация исходного кода, защищённых записей и подробностей инфраструктуры или безопасности не требуется; самостоятельного критерия либо исключения не вводится. Законные временные правила сбоя отмечаются в записи версии по пункту 11.

\subsection*{Глава 3. Качество данных и сбои}
8. В расчёте фиксируются источник, период, время получения и состояние проверки данных. Настоящий ноль отличается от неразрешённых данных. Дата либо период, описываемые исходной записью, отличаются от времени её получения. Необходимые входные значения либо законно сохранённый снимок источника позволяют восстановить расчёт с соблюдением пункта 28. Состояния «значение подтверждено», «данные не получены», «требуется обновление», «данные противоречивы» и «формат неверен» сохраняются раздельно. Для использованной записи фиксируются идентификатор источника и записи, относящийся к ней период, дата получения и применённое преобразование. Для конфликта сохраняются обе версии значения и основание выбора; более позднее получение само по себе не означает большей достоверности. Пользовательская поправка содержит автора, время и законное основание.

9. Не допускается автоматически приравнивать неразрешённые данные к ухудшающему положение семьи значению. При прямо предусмотренном законодательством нормативном вменении основание и причина показываются отдельно.

10. При невозможности расчёта из-за неразрешённых данных уполномоченный работник «Инсон» регистрирует вопрос в YAMIH и направляет ответственному собственнику либо оператору источника. Ошибка программы направляется системному подразделению; неопределённость правового правила рассматривается методическим подразделением и юридической службой и при необходимости официального толкования или изменения вносится компетентному органу. Работники Агентства не подменяют его полномочия. Ответственный работник отслеживает законный срок, сообщает семье результат и способ оспаривания, передаёт нерешённый вопрос руководителю до истечения срока. Состояние не является новым основанием отказа или прекращения. Документ семьи вместо данных государственной системы допускается только в предусмотренных законодательством случаях.

11. При законно разрешённом временном изменении алгоритма из-за технической неисправности, перерыва обмена либо неполного получения электронных данных уполномоченное должностное лицо фиксирует основание, область, начало, временное правило и условия восстановления. Такие расчёты отдельно отмечаются; после восстановления проверка и необходимый повторный расчёт выполняются по законодательству. Временная запись связывается с затронутыми источниками, расчётами и версиями; указываются принявший решение работник, правовая основа его полномочия и дата прекращения временного режима. После восстановления сопоставляются временный и восстановленный расчёты, сохраняются расхождения и индивидуальные решения компетентного органа. Закрытие сбоя подтверждается доступностью источника и завершением требуемых проверок, а не стиранием временной записи.

\subsection*{Глава 4. Проверка расчётов}
12. До внедрения либо изменения, способного повлиять на балл, категорию или основание решения, реализация сравнивается с независимо обоснованными контрольными ответами из утверждённого правового правила. Существующие приёмочные проверки и проверки технического соответствия включают применимые границы, исключения, их взаимодействия, состояния качества данных и даты действия. Документированный эквивалентный тест используется без повторения только ради настоящего Порядка. Ожидаемый ответ не получается исключительно из проверяемой реализации. Внутренний регламент определяет реквизиты и сверку запланированных и исполненных случаев. Пропуск применимого пути, пустой состав сравнения и неразрешённый конфликт правила не представляются успешной проверкой.

13. До получения результатов задаются цель, период, случаи, отбор, сравнение и условия оценки. Последующие изменения записываются с причиной. Синтетические результаты не доказывают правильность определения реальной нуждаемости.

14. Несоответствия балла и категории учитываются отдельно. Указываются изменившиеся категории, соответствующее общее число и содержание изменения. Неясные случаи не скрываются и отделяются от определённых категорий. Знаменатель доли расхождений категории включает только сопоставимые случаи с определённой категорией в обоих расчётах при одном утверждённом правиле и одинаковых входах. Знаменатель чувствительности включает только парные случаи с определённой категорией до и после заданного изменения. Для каждой доли указываются числитель, знаменатель, исключённые и неразрешённые случаи; при нулевом знаменателе доля считается неопределённой, а не нулевой.

15. Ошибки нуждаемости рассчитываются только при независимо обоснованных данных сравнения. Прежнее автоматическое решение, старая категория или результат жалобы не являются сами по себе правильной оценкой. Жалобы не толкуются как показатель ошибок, охватывающий не подавших жалоб семьи.

16. Протокол указывает утверждённое правило и проверяемую реализацию, область и основание проверки, исполнителей, результаты, непроверенные и неразрешённые вопросы, выявленные последствия и решения об исправлении. Соответствие правилу, качество данных и полнота оснований отделяются от эмпирического определения нуждаемости. Расхождения промежуточных компонентов сохраняются и при совпадении итога. Совпадение программ само по себе не устанавливает законный ожидаемый ответ. Внутренний технический регламент определяет доказательства и формы записей; защищённые входы сохраняются в законно доступной части.

17. Несоответствующая методике версия исправляется и повторно проверяется. Затронутые случаи определяются и законно пересматриваются компетентным органом. Непроверенная версия не используется как подтверждённо соответствующая; это не основание остановки обращений или законной помощи. Для дефекта регистрируются правило и версия, воспроизводимый случай, затронутые решения, ответственный исполнитель, способ исправления, применимый срок и результат повторной проверки. Отбор затронутых решений основывается на использованных версии, датах и условиях, а не только на выявившей дефект жалобе. Изменённые компоненты и зависимые расчётные пути проверяются повторно; при невозможности определить область влияния проверяется вся применимая расчётная цепочка. Закрытие подтверждается результатом проверки и записью о требуемом пересмотре индивидуальных решений.

\subsection*{Глава 5. Основания и пересмотр решений}
18. Вместе с результатом сохраняются версия, критерии, источники и периоды, балл, правило категории, нормативное вменение и основание исправления работником. При повторении прежняя запись не стирается. Сохраняются необходимые промежуточные вклады, основания включения либо исключения, этапы преобразования и округления и идентификаторы пункта 6. Взаимно компенсирующие расхождения компонентов не скрываются совпадением итоговой суммы.

19. Вместе с решением через личный кабинет предоставляются правовые и фактические основания, версия, основные значения и периоды, причина вменения или исключения и способ возражения. Не использующий кабинет получает доступ через «Инсон». Защищённые данные третьих лиц не раскрываются. SMS не заменяет объяснение и доступ. Объяснение включает идентификатор и дату решения, применённое правовое правило, относящиеся к заявителю значения и периоды, основные преобразования, категорию и условия конкретной помощи либо направления, компетентного получателя возражения и законный срок. Нормативное вменение отличают от подтверждённого фактического дохода. Для решения только на основе автоматизированной обработки заранее фиксируются допустимое по закону основание и порядок разъяснения этой обработки и возможных правовых последствий. Ограничение раскрытия конкретных сведений имеет правовое основание и не исключает предоставления остальных доступных оснований.

20. Возражение против решения только на основе автоматизированной обработки рассматривает собственник либо оператор, письменно сообщающий результат субъекту в течение десяти дней. Возражение может направляться через YAMIH или «Инсон»; регистрация в системе не откладывает исчисление срока. Исправление данных направляется уполномоченному источнику. По обращению субъекта собственник или оператор изменяет либо дополняет персональные данные не позднее трёх дней с момента обращения; установленное несоответствие действительности исправляется немедленно. Жалоба совету рассматривается в его законной предметной компетенции. Неподведомственное вычислительное возражение по этой причине не оставляется без рассмотрения, а направляется компетентному органу. Сохраняются законные судебные и иные способы и сроки. Нового обязательного досудебного этапа не вводится.

\subsection*{Глава 6. Исследование, отчёт и контроль}
21. До научного использования определяются законное основание обработки, цель, необходимый объём и защита. Научное обезличивание соответствует законодательству. Замена идентификатора сама по себе не является обезличиванием. Соблюдаются специальные требования к здоровью. В записи исследовательского допуска указываются законный вопрос, необходимые поля и период, уполномоченный оператор, разрешённые действия, срок хранения и контроль раскрытия. Рабочее участие гражданина само по себе не является согласием семьи на обработку её данных. Информация о здоровье используется только при соблюдении специальных законных условий. Исследовательское вычисление не меняет категорию или выплату без отдельного законного индивидуального решения.

22. Исследовательский расчёт выполняется в среде Агентства. Порядок не даёт внешнему участнику права получить персональные данные. Защищённые записи, идентифицирующие малые группы и связываемые таблицы не публикуются. До исполнения предоставленного кода уполномоченный оператор проверяет его состав, версию, требуемый доступ и исходящие соединения. Обработка выполняется с минимально необходимыми правами; передача защищённых входов внешним сервисам и включение их в открытые журналы не допускаются. Разрешение на запуск и результаты контроля раскрытия фиксируются в протоколе. Предоставление кода не означает разрешения его автоматического запуска в производственной системе.

23. Протокол проверки хранится внутри Агентства совместно с существующими записями выпуска и контроля. Он указывает правило и программную версию, дату, область, исполненные и исключённые случаи, сопоставимые и неразрешённые количества, доказанные расхождения, меры исправления и пределы. На эквивалентные записи можно ссылаться без создания повторного журнала. Раскрытие выполняется через существующие законные процедуры отчётности, доступа заявителя и контроля. Отдельного отчёта перед каждым выпуском, новой периодической статистической публикации и нового сбора данных семей не устанавливается. Открытая выписка, когда законно обязательна или разрешена, проверяется на раскрытие, указывает область и пределы; защищённые и связываемые записи семей исключаются.

24. Контроль хранения версий, результатов и исправлений выполняется в пределах действующих полномочий. Проверка служит совершенствованию методики, не самостоятельному основанию автоматического сокращения помощи.

\subsection*{Глава 7. Приёмка, исправление и хранение}
25. Руководитель Агентства или законно уполномоченный заместитель подтверждает соответствие расчётов в существующей записи выпуска с методическим, системным и юридическим заключениями. Запись связывает требуемые законодательством согласования, правовые даты применения, результаты независимых контролей, используемые эквивалентные проверки, дополнительные работы и ресурсы. Открытые вопросы и действия по пункту 26 фиксируются. Внутреннее подтверждение не заменяет компетентное техническое согласование и согласование кибербезопасности, не даёт полномочия определять критерии нуждаемости или официальное толкование. Доказанное влияющее на решение расхождение не считается несущественной оставшейся задачей.

26. Выпуск с доказанным расхождением балла, категории либо основания решения при одном утверждённом правиле, одинаковых входах и утверждённом округлении исправляется; затронутые пути проверяются повторно до их внедрения. Незатронутые пути продолжаются только при доказанном отделении и соответствии. Предыдущий программный выпуск сохраняется или восстанавливается только при исполнении законно применимого на дату решения правила и проверке относящихся к нему контролей и пригодности к работе; возраст выпуска не продлевает отменённую методику. При выявлении дефекта в работе системное подразделение фиксирует область и затронутые заявления, приоритетно исправляет дефект и направляет индивидуальные решения для законного пересмотра. Даты поступления, законные сроки и компетентная процедура сохраняются; техническое ожидание не является новым основанием отказа или полномочием остановить законную помощь. Полномочие при сбое обмена действует только в своей законной предметной области. Не даётся права ручного внесения категории, применения заменённых правил оценки, освобождения от долга или приостановления январского внедрения. Отсутствующее либо противоречивое правовое правило Агентство и Министерство вносят компетентному принимающему или толкующему органу для разрешения до первого применения; внутренняя настройка и толкование не изменяют право. Технический резерв не решает отсутствие законного правила. Конечные контроли проверяют указанную область, не все входы и не реальную точность нуждаемости.

27. Новые критерии действуют с даты компетентного акта. Ошибка прежнего решения проверяется по применённому к нему праву; новый критерий и позднейший параметр автоматически назад не применяются. Повторение само по себе не основание взыскать ранее законно назначенную выплату. Изменение, продолжение и взыскание решаются по соответствующему основанию и индивидуальному решению. Запись пересмотра указывает, исправлены ли исходные сведения, исполнение прежнего правила или само правовое правило; сохраняются прежние и исправленные компоненты, применимая дата и основания изменения индивидуального решения. Перечень затронутых дел и результат их передачи компетентному органу связываются с записью дефекта по пункту 17. При отсутствии законного основания изменить выплату техническое исправление не превращается в платёжное поручение.

28. Агентство определяет категории записей, законный срок и основание хранения, доступ, историю исправлений, обезличивание или уничтожение. Порядок не устанавливает бессрочного хранения. При истечении основания или срока данные уничтожаются либо в законно разрешённом случае обезличиваются по архивным и персональным требованиям. Когда уничтожение обязательно, обезличивание его не заменяет. Необходимые для жалобы или контроля записи защищаются по действующим правилам хранения. Категории хранения охватывают исходные значения, расчётные следы, контрольные протоколы, записи исправлений и обращения раздельно. Для каждой указываются цель, правовое основание срока, ответственный, допустимые роли доступа и способ прекращения хранения. Доступ и изменение защищённых записей регистрируются; проверяющим предоставляется необходимый объём в пределах полномочий. Наличие контрольной копии не создаёт права хранить персональные сведения после прекращения законного основания.
\clearpage
\section*{Краткая альтернативная редакция для основного проекта}
Рекомендуемый минимум для нормативного включения в акт по пункту 14 PF-193. Пять пунктов ниже заменяют полную редакцию из семи основных пунктов и 28 пунктов приложения, а не добавляются к ней. Детальные паспорта, случаи и протоколы передаются во внутреннюю техническую документацию. Если конкретное требование уже обеспечено применимой нормой надлежащей юридической силы, согласованной с вышестоящими актами, разработчик сохраняет эту норму и указывает её в основном проекте либо его обосновании; повторение текста не требуется. Последующая нумерация определяется местом включения.

1. Национальному агентству социальной защиты при Президенте Республики Узбекистан (далее — Агентство) совместно с Министерством экономики и финансов Республики Узбекистан (далее — Министерство) при подготовке методики многомерной оценки для утверждения в установленном порядке обеспечить правовое определение критериев и условий распределения семей по категориям. Определение охватывает влияющие на категорию пороги, равенство, приоритет исключений, единицы и периоды, нормативное вменение, допустимые входные значения и округление в той части, где они устанавливают условие признания или отказа. Указываются компетентный акт, начало применения и отсылки к сохранённым действующим правилам. Описание программы и техническая приёмка не устанавливают отсутствующее правовое условие и не изменяют вышестоящую норму. Нормативные положения утверждаются и официально публикуются в установленном порядке; для ведомственного нормативного акта соблюдается применимое требование государственной регистрации. Прямо предусмотренные законодательством полномочия временного изменения алгоритма сохраняются в своей предметной области.

2. Агентству обеспечить соответствие расчёта законно применимой методике и проверку затронутых правил до первого применения либо относящегося к ним программного изменения в существующем процессе выпуска. Равноценные действующие проверки используются без дублирования. Результат связывается с правовым правилом, датами, проверяемым выпуском и независимо обоснованным ожидаемым ответом. Доказанное влияющее на решение несоответствие исправляется и проверяется повторно; другая сборка применяется только при её соответствии правилу, законно действующему на дату решения, и подтверждённой пригодности. Затронутые индивидуальные решения передаются компетентному органу для пересмотра в законном порядке. Сохраняются обязательные технические согласования, законные сроки, дата начала внедрения и условия назначения помощи.

3. Вместе с юридически значимым результатом Агентству сохранять законно необходимую запись применённого правового правила, версии расчёта, исходных значений и периодов, основных преобразований, вменения, исключения и основания категории, позволяющую восстановить расчёт. Через действующие личный кабинет либо центр социальных услуг «Инсон» Агентства заявителю предоставляются доступные правовые и фактические основания, относящиеся к нему значения и периоды и порядок возражения; SMS результата не заменяет объяснение. Сохраняются законные сроки исправления данных и рассмотрения возражений. Защищённые сведения третьих лиц не раскрываются; хранение, доступ и прекращение хранения соответствуют законодательству.

4. Изменение правового условия оценки вносится компетентному органу и применяется по законно установленной дате; исправление программы исполняет неизменное законное правило. Если законодательство прямо разрешает временное изменение алгоритма при технических неисправностях, перерывах обмена либо неполном получении электронных данных, уполномоченное лицо фиксирует основание, затронутые источники и случаи, временное правило, начало и условия прекращения, а после восстановления — необходимые проверки и законный пересмотр. Такое полномочие не расширяется настоящим предложением. Отсутствие либо противоречие правовой нормы разрешается компетентным принимающим или толкующим органом до первого применения затронутого условия; толкование не вносит поправок. Техническое ожидание не создаёт основания отказа, продления срока, применения отменённого критерия или освобождения от долга.

5. Проверки, формы и записи организуются Агентством в пределах действующих полномочий и существующих процессов разработки, согласования и контроля. Внутренние технические документы не устанавливают новых обязанностей заявителя или правовых критериев; эквивалентные спецификации, тесты, записи и каналы используются повторно. До обязательства дополнительных расходов Агентство совместно с Министерством определяет только непокрытую работу, назначенное время, договорное покрытие и законное финансирование; изменение обработки заявлений обосновывается относящимися к нему фактическими нагрузкой и временем. Не создаются отдельная ручная служба, дополнительный этап заявления или новая обязательная публичная статистическая отчётность.

\end{document}
